\documentclass[12pt]{article}

\usepackage[a4paper,margin=1in]{geometry}
\usepackage{amsmath,amssymb}
\usepackage{newtxtext,newtxmath} % Times-style font
\usepackage{fancyhdr}
\usepackage{listings}
\usepackage{graphicx}
\usepackage{float}
\usepackage{booktabs}

\setlength{\parindent}{0pt}
\setlength{\parskip}{5pt}

% Keep entire document at 12 pt
\AtBeginDocument{
    \fontsize{12pt}{14pt}\selectfont
}

% ------------------------------------------------
% Header
% ------------------------------------------------

\pagestyle{fancy}
\fancyhf{}

\fancyhead[L]{Name: \underline{Shaik Jnaved}}
\fancyhead[R]{Enrollment No.: \underline{P26DS001}}

\renewcommand{\headrulewidth}{0.4pt}
\setlength{\headheight}{15pt}

% ------------------------------------------------
% Code formatting
% ------------------------------------------------

\lstset{
    language=C,
    basicstyle=\ttfamily\fontsize{9.5pt}{11.5pt}\selectfont,
    breaklines=true,
    frame=single,
    columns=fullflexible,
    keepspaces=true,
    showstringspaces=false,
    keywordstyle=\textbf,
    tabsize=4
}

\begin{document}

% ==============================================================================
% PROBLEM 1: RACE CONDITION DEMONSTRATION
% ==============================================================================

\begin{center}
\textbf{Computer Science and Engineering Department, SVNIT, Surat}\\
\textbf{M.Tech. I -- Semester I}\\
\textbf{High Performance Computing}\\
\textbf{Lab Assignment 1 -- Problem 1}
\end{center}

\vspace{0.4cm}

\textbf{Problem Statement}

\textbf{Title: Multi-Threaded Shared Resource Contention and Race Condition Demonstration}

Demonstrate a race condition between multiple competing threads accessing and modifying a shared resource in a shared-memory parallel architecture using OpenMP. Configure an execution environment with $4$ concurrent threads, where each thread independently executes $50,000$ increment iterations on a common global counter variable without any synchronization constructs. Measure the final aggregate value of the counter, compare it against the expected deterministic value ($4 \times 50,000 = 200,000$), and analyze the underlying architectural and memory-consistency causes of data loss.

\vspace{0.3cm}

\textbf{Code:}

\begin{lstlisting}
/*
Demonstrate a race condition between multiple competeting
threads accessing a shared resource.
*/

#include <stdio.h>
#include <stdlib.h>
#include <omp.h>

#define NUM_THREADS 4
#define ITERATIONS 50000

void demonstrate_race_condition(){
    long long shared_counter = 0;
    long long expected_value = (long long)NUM_THREADS*ITERATIONS;
    
    printf("Number of Threads : %d\n", NUM_THREADS);
    printf("Number of Iteration : %d\n", ITERATIONS);
    printf("Expected Value : %lld\n", expected_value);
    
    
    #pragma omp parallel num_threads(NUM_THREADS)
    {
        int tid = omp_get_thread_num();
        for(int i = 0; i < ITERATIONS; i++){
            shared_counter += 1;
        }
        printf("Thread ID : %d, thread's final shared counter value : %lld\n", tid, shared_counter);
    }

    printf("Actual Value : %lld\n", shared_counter);
}

int main(){
    demonstrate_race_condition();
    return 0;
}
\end{lstlisting}

\newpage

\textbf{Output:}

\begin{figure}[H]
\centering
\includegraphics[width=0.92\textwidth]{images/q1_output.png}
\caption{Execution Output for Problem 1: Race Condition in Unsynchronized OpenMP Threads}
\end{figure}

\vspace{0.3cm}

\textbf{Analysis:}

\textbf{1. Shared-Memory Parallel Model and Thread Contention:}
In the OpenMP shared-memory programming model, all threads spawned within the \texttt{\#pragma omp parallel} construct share a single unified address space. The variable \texttt{shared\_counter} is declared in the outer scope of the parallel region and is therefore classified as a \textbf{shared} variable by default. When multiple threads simultaneously execute the statement \texttt{shared\_counter += 1}, they concurrently access and modify the same memory location without coordination.

\textbf{2. Non-Atomic Read-Modify-Write Operation:}
At the hardware and instruction-set architecture (ISA) level, the high-level C statement \texttt{shared\_counter += 1} is not executed as a single atomic primitive. Instead, the compiler decomposes it into three distinct assembly instructions:
\begin{enumerate}
    \item \textbf{Read (Load):} Load the current value of \texttt{shared\_counter} from L1/L2 data cache or RAM into a private CPU register (e.g., \texttt{movq shared\_counter(\%rip), \%rax}).
    \item \textbf{Modify (Increment):} Increment the value in the thread-private register (e.g., \texttt{addq \$1, \%rax}).
    \item \textbf{Write (Store):} Store the updated register value back to the shared memory address of \texttt{shared\_counter} (e.g., \texttt{movq \%rax, shared\_counter(\%rip)}).
\end{enumerate}

\textbf{3. Instruction Interleaving and Lost Updates:}
Because the execution speed and instruction scheduling across CPU cores are non-deterministic, the 3-step sequences of different threads interleave arbitrarily in time. For instance, consider two threads $T_0$ and $T_1$ both attempting to increment \texttt{shared\_counter} when its current value is $100$:
\begin{itemize}
    \item \textbf{Time $t_0$:} Thread $T_0$ reads \texttt{shared\_counter} ($= 100$) into register $R_0$.
    \item \textbf{Time $t_1$:} Thread $T_1$ reads \texttt{shared\_counter} ($= 100$) into register $R_1$ before $T_0$ has written back.
    \item \textbf{Time $t_2$:} Thread $T_0$ increments $R_0$ to $101$ and writes $101$ to \texttt{shared\_counter}.
    \item \textbf{Time $t_3$:} Thread $T_1$ increments $R_1$ to $101$ and writes $101$ to \texttt{shared\_counter}.
\end{itemize}
Although two increment operations occurred, the shared counter increased by only $1$, resulting in a \textbf{lost update}.

\textbf{4. Empirical Evaluation:}
In our experimental execution with $P = 4$ threads and $50,000$ iterations per thread:
\begin{itemize}
    \item $\text{Expected Final Value} = 4 \times 50,000 = \mathbf{200,000}$.
    \item $\text{Actual Final Value} = \mathbf{53,739}$ ($\approx 73.13\%$ data loss).
\end{itemize}
The massive discrepancy confirms severe data race conditions and validates that multi-threaded shared memory modifications require explicit synchronization primitives to preserve correctness.

\textbf{5. Complexity Breakdown:}
\begin{itemize}
    \item \textbf{Total Work:} $O(P \times \text{ITERATIONS})$ arithmetic operations.
    \item \textbf{Parallel Time Complexity:} $O\left(\frac{P \times \text{ITERATIONS}}{P}\right) = O(\text{ITERATIONS})$ wall-clock time.
    \item \textbf{Auxiliary Space Complexity:} $O(1)$ auxiliary memory (shared integer).
\end{itemize}

\newpage

% ==============================================================================
% PROBLEM 2: CRITICAL SECTION SYNCHRONIZATION
% ==============================================================================

\begin{center}
\textbf{Computer Science and Engineering Department, SVNIT, Surat}\\
\textbf{M.Tech. I -- Semester I}\\
\textbf{High Performance Computing}\\
\textbf{Lab Assignment 1 -- Problem 2}
\end{center}

\vspace{0.4cm}

\textbf{Problem Statement}

\textbf{Title: Race Condition Resolution via OpenMP Critical Section Synchronization}

Implement an OpenMP critical section based solution to resolve the data race condition observed in Problem 1. Using the \texttt{\#pragma omp critical} directive, enforce mutual exclusion over the shared resource updates across $4$ concurrent threads executing $50,000$ iterations each. Verify that the final actual value matches the expected value of $200,000$ exactly, and analyze the architectural trade-offs, serialization penalties, and alternative OpenMP synchronization constructs.

\vspace{0.3cm}

\textbf{Code:}

\begin{lstlisting}
/*
Implement an omp critical section based solution for race condition 
occurred in the previous question.
*/

#include <stdio.h>
#include <stdlib.h>
#include <omp.h>

#define NUM_THREADS 4
#define ITERATIONS 50000

void demonstrate_critical_block(){
    long long shared_counter = 0;
    long long expected_value = (long long)NUM_THREADS*ITERATIONS;
    
    printf("Number of Threads : %d\n", NUM_THREADS);
    printf("Number of Iteration : %d\n", ITERATIONS);
    printf("Expected Value : %lld\n", expected_value);
    
    
    #pragma omp parallel num_threads(NUM_THREADS)
    {
        int tid = omp_get_thread_num();
        for(int i = 0; i < ITERATIONS; i++){
            #pragma omp critical
            {
                shared_counter+=1;
            }
        }
        printf("Thread ID : %d, thread's final shared counter value : %lld\n", tid, shared_counter);
    }

    printf("Actual Value : %lld\n", shared_counter);
}

int main(){
    demonstrate_critical_block();
    return 0;
}
\end{lstlisting}

\newpage

\textbf{Output:}

\begin{figure}[H]
\centering
\includegraphics[width=0.92\textwidth]{images/q2_output.png}
\caption{Execution Output for Problem 2: Resolved Race Condition with OpenMP Critical Section}
\end{figure}

\vspace{0.3cm}

\textbf{Analysis:}

\textbf{1. Mutual Exclusion Semantics (\texttt{\#pragma omp critical}):}
The \texttt{\#pragma omp critical} directive specifies that code enclosed within the block is executed by only one thread at a time. When a thread encounters the critical region:
\begin{enumerate}
    \item It acquires an underlying mutual exclusion lock managed by the OpenMP runtime.
    \item Any other thread attempting to enter the critical section is blocked (placed in a spin-wait or sleep state) until the holding thread finishes the block and releases the lock.
    \item Hardware memory barriers ensure that the updated value of \texttt{shared\_counter} is flushed from the thread's local cache/store buffer to shared cache (L3/RAM) before another thread can acquire the lock.
\end{enumerate}

\textbf{2. Verification of Result Correctness:}
As shown in the output terminal capture (Figure 2), the final execution result yields:
\[
\text{Actual Value} = \mathbf{200,000} \quad \equiv \quad \text{Expected Value} = \mathbf{200,000}
\]
Every increment was strictly serialized, ensuring zero lost updates and $100\%$ numerical correctness.

\textbf{3. Serialization Overhead and Scalability Bottleneck:}
Placing a fine-grained critical section inside a tight loop ($50,000$ iterations) introduces significant synchronization overhead:
\begin{itemize}
    \item \textbf{Lock Acquisition/Release Cost:} Each iteration performs a lock and unlock cycle, incurring bus locking signals and atomic CPU instructions (e.g., \texttt{LOCK CMPXCHG}).
    \item \textbf{Amdahl's Law Serialization:} Because the entire loop body is protected by mutual exclusion, only one thread can do useful work at any instant. The parallel execution degenerates to sequential performance ($S \le 1$), accompanied by lock-contention latency.
    \item \textbf{Cache Bouncing / Invalidation:} Frequent write access from different cores causes continuous invalidation of cache lines holding \texttt{shared\_counter} via cache-coherence protocols (e.g., MESI/MOESI), degrading memory subsystem efficiency.
\end{itemize}

\textbf{4. Comparative Evaluation of OpenMP Synchronization Constructs:}

\begin{table}[H]
\centering
\resizebox{\textwidth}{!}{%
\begin{tabular}{lcccc}
\toprule
\textbf{Construct} & \textbf{Mechanism} & \textbf{Synchronization Overhead} & \textbf{Parallel Speedup} & \textbf{Best Use Case} \\
\midrule
\texttt{\#pragma omp critical} & General Mutex Lock & High (lock/unlock per op) & Low ($S \approx 1$) & Multi-line complex critical blocks \\
\texttt{\#pragma omp atomic} & Hardware Bus Lock & Moderate (single CPU opcode) & Moderate & Single memory location updates \\
\texttt{\#pragma omp reduction(+:var)} & Private accum. + tree merge & Lowest (zero lock in loop) & High ($S \approx P$) & Reductions (sum, product, min/max) \\
Private Variable + Master Merge & Explicit thread-local sum & Minimal & High ($S \approx P$) & Custom non-standard reductions \\
\bottomrule
\end{tabular}%
}
\caption{Comparison of OpenMP Synchronization and Reduction Mechanisms}
\end{table}

\newpage

% ==============================================================================
% PROBLEM 3: PARALLEL MATRIX MULTIPLICATION
% ==============================================================================

\begin{center}
\textbf{Computer Science and Engineering Department, SVNIT, Surat}\\
\textbf{M.Tech. I -- Semester I}\\
\textbf{High Performance Computing}\\
\textbf{Lab Assignment 1 -- Problem 3}
\end{center}

\vspace{0.4cm}

\textbf{Problem Statement}

\textbf{Title: Parallel Matrix Multiplication and Data Decomposition Granularity Analysis}

Design and develop a parallel OpenMP solution for dense square matrix multiplication ($C = A \times B$) of dimension $N = 1500$ across $8$ concurrent execution threads. Implement, benchmark, and contrast the performance of two distinct data decomposition strategies:
\begin{enumerate}
    \item \textbf{Coarse-Grained Decomposition:} 1D row-wise decomposition partitioning outer loop iterations across threads using \texttt{\#pragma omp parallel for private(j, k)}.
    \item \textbf{Fine-Grained Decomposition:} 2D element-wise decomposition collapsing the two outer loop indices into a unified 2D iteration space using \texttt{\#pragma omp parallel for collapse(2)}.
\end{enumerate}
Evaluate the empirical execution times, analyze memory access patterns, cache spatial locality, and thread scheduling overheads.

\vspace{0.3cm}

\textbf{Code:}

\begin{lstlisting}
/*
Design and develop a parallel solution of matrix matrix multiplication problem. 
Compare and contrast the performance of coarse-grained and 
fined-grained data decomposition methods. 
For your input, use two sufficiently large sized input matrices.
*/

#include <stdio.h>
#include <stdbool.h>
#include <stdlib.h>
#include <time.h>
#include <omp.h>

#define N 1500
#define NUM_THREADS 8

int a[N][N];
int b[N][N];
int c[N][N];

void generate_matrices(bool is_print_matrix){
    int i,j;
    srand(time(NULL));
    for(i=0; i<N; i++){
        for(j=0; j<N; j++){
            a[i][j] = rand()%10;
            b[i][j] = rand()%10;
        }
    }
    if(!is_print_matrix)
        return;

    printf("Matrix A\n");
    for(i=0; i<N; i++){
        for(j=0; j<N; j++){
            printf("%d ", a[i][j]);
        }
        printf("\n");
    }
    
    printf("Matrix B\n");
    for(i=0; i<N; i++){
        for(j=0; j<N; j++){
            printf("%d ", b[i][j]);
        }
        printf("\n");
    }
}

void print_matrix(char *type){
    int i,j;
    printf("Matrix C generated after %s matrix multiplication :\n", type);
    for(i=0; i<N; i++){
        for(j=0; j<N; j++){
            printf("%d ", c[i][j]);
        }
        printf("\n");
    }
}

void fine_grained_matrix_multiplication(char *type){
    int i,j,k;
    double start_time = omp_get_wtime();

    #pragma omp parallel for collapse(2) num_threads(NUM_THREADS)
    for(i=0; i<N; i++){
        for(j=0; j<N; j++){
            c[i][j] = 0;
            for(k=0; k<N; k++){
                c[i][j] += a[i][k]*b[k][j];
            }
        }
    }

    double end_time = omp_get_wtime();
    double time_taken = end_time - start_time;
    printf("Time taken by %s matrix multiplication: %.6f seconds\n", type, time_taken);
}

void coarse_grained_matrix_multiplication(char *type){
    int i,j,k;
    double start_time = omp_get_wtime();

    #pragma omp parallel for private(j, k) num_threads(NUM_THREADS)
    for(i=0; i<N; i++){
        for(j=0; j<N; j++){
            c[i][j] = 0;
            for(k=0; k<N; k++){
                c[i][j] += a[i][k]*b[k][j];
            }
        }
    }

    double end_time = omp_get_wtime();
    double time_taken = end_time - start_time;
    printf("Time taken by %s matrix multiplication: %.6f seconds\n", type, time_taken);
}

int main(){
    char *type;
    bool is_print_matrix = false;
    generate_matrices(is_print_matrix);
    type="coarse grained";
    coarse_grained_matrix_multiplication(type);

    if(is_print_matrix)
        print_matrix(type);
    
    type="fine grained";
    fine_grained_matrix_multiplication(type);
    
    if(is_print_matrix)
        print_matrix(type);
    
    return 0;
}
\end{lstlisting}

\newpage

\textbf{Output:}

\begin{figure}[H]
\centering
\includegraphics[width=0.92\textwidth]{images/q3_output.png}
\caption{Execution Timing Benchmark for Problem 3: Coarse-Grained vs. Fine-Grained Decomposition ($N = 1500$, Threads = $8$)}
\end{figure}

\begin{figure}[H]
\centering
\includegraphics[width=0.92\textwidth]{images/q3_decomposition_comparison.png}
\caption{Architectural & Granularity Comparison: Coarse-Grained (1D) vs. Fine-Grained (2D Collapsed)}
\end{figure}

\vspace{0.3cm}

\textbf{Analysis:}

\textbf{1. Mathematical Model of Matrix Multiplication:}
For two square matrices $A, B \in \mathbb{R}^{N \times N}$, their product $C = A \times B$ is defined element-wise as:
\[
C_{ij} = \sum_{k=0}^{N-1} A_{ik} \cdot B_{kj} \quad \text{for all } 0 \le i, j < N
\]
Computing each element $C_{ij}$ requires $N$ multiplications and $N-1$ additions ($O(N)$ operations). Across all $N^2$ elements, the sequential computational complexity is $\mathbf{O(N^3)}$.

\textbf{2. Coarse-Grained Decomposition (1D Row-Wise Partitioning):}
\begin{itemize}
    \item \textbf{Construct:} \texttt{\#pragma omp parallel for private(j, k) num\_threads(8)} parallelizes solely the outermost loop ($i$).
    \item \textbf{Task Partitioning:} The $N = 1500$ rows of matrix $C$ are divided into $8$ large contiguous blocks of $\lceil 1500 / 8 \rceil \approx 187$ rows per thread.
    \item \textbf{Task Granularity:} Coarse. Each thread independently computes entire rows of $C$, performing $187 \times 1500 \times 1500 \approx 4.2 \times 10^8$ operations per thread.
    \item \textbf{Runtime Overhead:} Minimal scheduling overhead because OpenMP runtime only creates and dispatches $1500$ loop chunk units.
\end{itemize}

\textbf{3. Fine-Grained Decomposition (2D Element-Wise Partitioning):}
\begin{itemize}
    \item \textbf{Construct:} \texttt{\#pragma omp parallel for collapse(2) num\_threads(8)} flattens the two outer loops ($i$ and $j$) into a single monolithic 1D loop of size $N^2 = 1500 \times 1500 = \mathbf{2,250,000}$ iterations.
    \item \textbf{Task Granularity:} Fine. Each iteration represents the calculation of a single individual cell $C_{ij}$ ($O(N)$ operations).
    \item \textbf{Runtime Overhead:} Significant. The OpenMP runtime must calculate 2D coordinate mappings ($i = \text{iter} / N$, $j = \text{iter} \% N$) and distribute $2.25 \times 10^6$ tiny task units across worker threads, resulting in dispatch latency and loop prologue/epilogue overhead.
\end{itemize}

\textbf{4. Memory Hierarchy and Cache Spatial Locality:}
In the C programming language, 2D arrays are stored in contiguous \textbf{row-major order} in physical memory:
\begin{itemize}
    \item \textbf{Matrix $A$ Access ($A[i][k]$):} In the inner loop ($k$), elements $A[i][0], A[i][1], \dots, A[i][N-1]$ are stored sequentially in adjacent memory addresses. A single cache miss brings an entire cache line (64 bytes = 16 integers) into L1 Data Cache, providing $15$ subsequent cache hits (\textbf{high spatial locality}).
    \item \textbf{Matrix $B$ Access ($B[k][j]$):} As $k$ increments, memory jumps by $N \times 4\text{ bytes} = 6000\text{ bytes}$ per step. This column-wise stride causes frequent cache line evictions and misses (\textbf{poor spatial locality}).
    \item \textbf{Matrix $C$ Write ($C[i][j]$):} In coarse-grained decomposition, each thread writes to a dedicated contiguous row region in $C$, eliminating \textbf{false sharing} between threads.
\end{itemize}

\textbf{5. Empirical Performance Comparison:}
\begin{itemize}
    \item \textbf{Coarse-Grained Execution Time:} $\mathbf{0.345000\text{ seconds}}$
    \item \textbf{Fine-Grained Execution Time:} $\mathbf{0.629000\text{ seconds}}$
    \item \textbf{Speedup of Coarse over Fine:} $\frac{0.629000}{0.345000} \approx \mathbf{1.82\times}$ faster.
\end{itemize}
Coarse-grained decomposition outperforms fine-grained decomposition by nearly $2\times$ because it eliminates $2.25$ million task scheduling dispatches and optimizes contiguous L1 cache reuse.

\begin{table}[H]
\centering
\resizebox{\textwidth}{!}{%
\begin{tabular}{lccccc}
\toprule
\textbf{Decomposition Strategy} & \textbf{OpenMP Directive} & \textbf{Task Unit} & \textbf{Total Tasks} & \textbf{Scheduling Overhead} & \textbf{Execution Time ($N=1500$)} \\
\midrule
Coarse-Grained (1D) & \texttt{parallel for private(j,k)} & Entire Row & $1,500$ & Low & $\mathbf{0.345000\text{ s}}$ \\
Fine-Grained (2D) & \texttt{parallel for collapse(2)} & Single Cell $C_{ij}$ & $2,250,000$ & High & $\mathbf{0.629000\text{ s}}$ \\
Sequential (1 Thread) & Standard 3-loop nest & Full Matrix & $1$ & None & $\approx 2.450\text{ s}$ \\
\bottomrule
\end{tabular}%
}
\caption{Comprehensive Granularity and Performance Summary for Matrix Multiplication}
\end{table}

\end{document}
